\documentclass[10pt,a4paper]{amsart}
\usepackage[margin=19mm]{geometry}
\usepackage{amsmath,amssymb,mathtools}
\usepackage[scr=rsfs,cal=euler]{mathalfa}
\usepackage{xfrac}
\usepackage[unicode,hidelinks,bookmarks=false]{hyperref}
\newcommand{\ie}{i.e.\ }
\newcommand{\cf}{cf.\ }
\newcommand{\resp}{respectively\ }
\newcommand{\Subsection}{\subsection}
\providecommand{\nobreakdash}{\nobreak\hbox{-}}
\setlength{\emergencystretch}{2em}
\multlinegap0pt
\usepackage{amssymb}
\usepackage{float}
\newtheorem{definition}{Definition}
\newtheorem{theorem}{Theorem}
\title{A Riemannian Extension of a Quadrature Surface Free Boundary Problem: Stability and Optimality}
\author{Ababacar Sadikhe Djite, Diaraf Seck}
\date{Source: arXiv:2608.29462v1 (CC BY 4.0)}
\begin{document}
\maketitle

\noindent\emph{Source and licence.} A Riemannian Extension of a Quadrature Surface Free Boundary Problem: Stability and Optimality. Version arXiv:2608.29462v1 (CC BY 4.0), distributed by arXiv under the Creative Commons Attribution 4.0 International licence, http://creativecommons.org/licenses/by/4.0/, as recorded in the arXiv licence metadata for this exact version and shown on the abstract page https://arxiv.org/abs/2608.29462v1. Full licence notice: \texttt{licenses/arxiv-2608.29462-CC-BY-4.0.txt}.\par\smallskip

\noindent\textit{Editorial scope.} Counted unit: the article's theorem thm:compactness-admissible (source0.tex lines 1139--1161). The article's complete original proof of it is delivered with it (source0.tex lines 1163--1544, one \texttt{proof} environment of 382 lines). No mathematics is written, repaired or re-derived: the statement and the proof are the article's own lines, in the article's own notation, and no retained line is substituted or re-flowed.  Retained uncounted as the source-local context the proof needs: lines 710--713; lines 912--980.  Nothing was omitted from any retained span, so the package carries no omission record.  No retained line contains a citation, so the wrapper carries no bibliography and no bibliography entry.  No retained line contains a figure environment, an image-inclusion command or drawing code, so no image is delivered and the package declares no asset dependency.  Editorial formatting only, in addition: an AMS \texttt{amsart} preamble that restates the article's own declarations and macros used by the retained lines, namely the environments \texttt{definition}, \texttt{theorem} on separate counters, together with the standard packages those lines need.  The retained environments print in this document as Definition 1, Theorem 1 in order of appearance, whereas the article prints the same items with the numbering of its own full document.  The complete native source of the article as distributed in the arXiv e-print for this exact version is bundled unchanged as \texttt{sources/208854/source0.tex}.
\begin{equation}
K\subset C.
\label{eq:K-compact-C}
\end{equation}

\begin{definition}[Riemannian $RC$-GNP condition]
\label{def:RC-GNP}
Let $C\subset M$ be the fixed reference domain satisfying
\eqref{eq:K-compact-C}. We say that a domain $\Omega\subset M$
satisfies the \emph{Riemannian $RC$-GNP condition relative to $C$} if
the following properties hold:

\begin{enumerate}
\item
\begin{equation}
C\subset\Omega;
\label{eq:RCGNP-inclusion}
\end{equation}

\item
\[
\partial\Omega\in C^{2,\alpha};
\]

\item there exist constants
\[
A>0,
\qquad
c_0>0,
\qquad
\rho_0>0,
\]
independent of $\Omega$, and a function
\[
h_\Omega\in C^{2,\alpha}(M)
\]
such that
\begin{equation}
\Omega=\{x\in M:h_\Omega(x)>0\},
\qquad
\partial\Omega=\{x\in M:h_\Omega(x)=0\},
\label{eq:defining-function}
\end{equation}
with
\begin{equation}
\|h_\Omega\|_{C^{2,\alpha}(M)}
\leq A,
\label{eq:uniform-C2alpha}
\end{equation}
and
\begin{equation}
|\nabla_g h_\Omega|_g
\geq c_0
\qquad
\text{whenever }|h_\Omega|\leq\rho_0;
\label{eq:uniform-gradient}
\end{equation}

\item the normal exponential map of $\partial\Omega$ is uniformly
nondegenerate: there exists $\rho_1>0$, independent of $\Omega$, such
that
\begin{equation}
\mathcal E_\Omega:
\partial\Omega\times(-\rho_1,\rho_1)
\longrightarrow M,
\qquad
\mathcal E_\Omega(x,t)
=
\exp_x(t\nu_\Omega(x)),
\label{eq:normal-exp-Omega}
\end{equation}
is a diffeomorphism onto its image.
\end{enumerate}
\end{definition}

\begin{theorem}[Compactness of the admissible class]
\label{thm:compactness-admissible}
Let $(M,g)$ be a compact smooth Riemannian manifold without boundary,
and let $(\Omega_j)_{j\in\mathbb N}\subset\mathcal A_C(M)$.
Assume that the admissible class satisfies the uniform RC--GNP
conditions of Definition~\ref{def:RC-GNP}. Then, up to extraction of
a subsequence, there exists a domain
$\Omega\in\mathcal A_C(M)$ such that
\begin{equation}
d_H(\overline{\Omega_j},\overline{\Omega})\longrightarrow0,
\label{eq:compactness-Hausdorff}
\end{equation}
\begin{equation}
\Omega_j\overset{K}{\longrightarrow}\Omega,
\label{eq:compactness-K}
\end{equation}
and
\begin{equation}
\chi_{\Omega_j}\longrightarrow\chi_\Omega
\qquad\text{strongly in }L^1(M,dV_g).
\label{eq:compactness-L1}
\end{equation}
\end{theorem}

\begin{proof}
For every $j\in\mathbb N$, let
\[
h_j:=h_{\Omega_j}\in C^{2,\alpha}(M)
\]
be a defining function associated with $\Omega_j$, so that
\begin{equation}
\Omega_j=\{h_j>0\},
\qquad
\partial\Omega_j=\{h_j=0\}.
\label{eq:compactness-defining-functions}
\end{equation}
By the uniform RC--GNP assumption, there exists a constant
$A>0$, independent of $j$, such that
\begin{equation}
\|h_j\|_{C^{2,\alpha}(M)}\le A.
\label{eq:compactness-uniform-C2alpha}
\end{equation}

Since $M$ is compact, the embedding
\[
C^{2,\alpha}(M)\subset C^2(M)
\]
is compact. Hence, after extracting a subsequence, there exists
\[
h\in C^{2,\alpha}(M)
\]
such that
\begin{equation}
h_j\longrightarrow h
\qquad\text{strongly in }C^2(M).
\label{eq:compactness-C2-convergence}
\end{equation}
In particular,
\[
h_j\longrightarrow h
\quad\text{uniformly on }M,
\]
and
\[
\nabla_g h_j\longrightarrow\nabla_g h
\quad\text{uniformly on }M.
\]

Define
\begin{equation}
\Omega:=\{h>0\}.
\label{eq:compactness-limit-domain}
\end{equation}

We first prove that $0$ is a regular value of $h$. Let
$x\in M$ be such that
\[
h(x)=0.
\]
Since $h_j\to h$ uniformly, we have
\[
|h_j(x)|<\rho_0
\]
for all sufficiently large $j$, where $\rho_0>0$ is the constant
appearing in the uniform non-degeneracy condition of the RC--GNP
assumption. Therefore,
\[
|\nabla_g h_j(x)|_g\ge c_0
\]
for all sufficiently large $j$. Passing to the limit and using the
uniform convergence of $\nabla_g h_j$ gives
\[
|\nabla_g h(x)|_g\ge c_0.
\]
Consequently,
\begin{equation}
|\nabla_g h|_g\ge c_0
\qquad\text{on }\{h=0\}.
\label{eq:compactness-limit-nondegeneracy}
\end{equation}
Thus $0$ is a regular value of $h$, and therefore
\[
\partial\Omega=\{h=0\}
\]
is a $C^{2,\alpha}$ embedded hypersurface.

In fact, the same argument shows that
\begin{equation}
|\nabla_g h(x)|_g\ge c_0
\qquad
\text{whenever }|h(x)|<\rho_0.
\label{eq:compactness-limit-nondegeneracy-neighborhood}
\end{equation}

We next prove the convergence of the boundaries. We claim that
\begin{equation}
d_H(\partial\Omega_j,\partial\Omega)
\longrightarrow0.
\label{eq:compactness-boundary-Hausdorff}
\end{equation}

Suppose first, by contradiction, that
\[
\sup_{x\in\partial\Omega_j}
d_g(x,\partial\Omega)
\not\longrightarrow0.
\]
Then there exist $\varepsilon>0$, a subsequence, and points
$x_j\in\partial\Omega_j$ such that
\[
d_g(x_j,\partial\Omega)\ge\varepsilon.
\]
Since $M$ is compact, after passing to a further subsequence we may
assume that
\[
x_j\longrightarrow x\in M.
\]
Because $x_j\in\partial\Omega_j$, we have
\[
h_j(x_j)=0.
\]
Hence
\[
|h(x_j)|
=
|h(x_j)-h_j(x_j)|
\le
\|h-h_j\|_{L^\infty(M)}
\longrightarrow0.
\]
By continuity of $h$,
\[
h(x)=0.
\]
Thus
\[
x\in\partial\Omega.
\]
It follows that
\[
d_g(x_j,\partial\Omega)
\le d_g(x_j,x)
\longrightarrow0,
\]
which contradicts
\[
d_g(x_j,\partial\Omega)\ge\varepsilon.
\]
Therefore,
\begin{equation}
\sup_{x\in\partial\Omega_j}
d_g(x,\partial\Omega)
\longrightarrow0.
\label{eq:compactness-boundary-one-sided}
\end{equation}

Conversely, suppose that
\[
\sup_{x\in\partial\Omega}
d_g(x,\partial\Omega_j)
\not\longrightarrow0.
\]
Then there exist $\varepsilon>0$, a subsequence, and points
$x_j\in\partial\Omega$ such that
\[
d_g(x_j,\partial\Omega_j)\ge\varepsilon.
\]
After extraction, we may assume
\[
x_j\longrightarrow x\in\partial\Omega.
\]

Since $0$ is a regular value of $h$, there exists a neighborhood $U$
of $x$ in which $h$ takes both positive and negative values. Hence
there exist points
\[
y^+,y^-\in U
\]
such that
\[
h(y^+)>0,
\qquad
h(y^-)<0.
\]
By the uniform convergence $h_j\to h$, for all sufficiently large $j$,
\[
h_j(y^+)>0,
\qquad
h_j(y^-)<0.
\]
Let $\gamma:[0,1]\to U$ be a continuous curve joining $y^-$ to
$y^+$. Then
\[
h_j(\gamma(0))<0,
\qquad
h_j(\gamma(1))>0.
\]
By the intermediate value theorem, there exists
$t_j\in(0,1)$ such that
\[
h_j(\gamma(t_j))=0.
\]
Consequently,
\[
z_j:=\gamma(t_j)\in\partial\Omega_j.
\]
Since $U$ can be chosen arbitrarily small around $x$ and
$x_j\to x$, this contradicts
\[
d_g(x_j,\partial\Omega_j)\ge\varepsilon.
\]
Therefore,
\begin{equation}
\sup_{x\in\partial\Omega}
d_g(x,\partial\Omega_j)
\longrightarrow0.
\label{eq:compactness-boundary-other-sided}
\end{equation}
Combining \eqref{eq:compactness-boundary-one-sided} and
\eqref{eq:compactness-boundary-other-sided}, we obtain
\[
d_H(\partial\Omega_j,\partial\Omega)\longrightarrow0.
\]

We now prove the compact convergence of the domains. Let
\[
K_1\subset\Omega.
\]
Since $h>0$ on the compact set $K_1$, there exists
\[
m_1:=\min_{K_1}h>0.
\]
For sufficiently large $j$,
\[
\|h_j-h\|_{L^\infty(M)}<\frac{m_1}{2}.
\]
Consequently,
\[
h_j(x)
\ge h(x)-\frac{m_1}{2}
\ge\frac{m_1}{2}>0
\]
for every $x\in K_1$. Hence
\[
K_1\subset\Omega_j
\]
for all sufficiently large $j$.

Similarly, let
\[
K_2\subset M\setminus\overline{\Omega}.
\]
Since $h<0$ on $K_2$, we may set
\[
m_2:=-\max_{K_2}h>0.
\]
For sufficiently large $j$,
\[
\|h_j-h\|_{L^\infty(M)}<\frac{m_2}{2},
\]
and therefore
\[
h_j(x)
\le h(x)+\frac{m_2}{2}
\le-\frac{m_2}{2}<0
\]
for every $x\in K_2$. Hence
\[
K_2\subset M\setminus\Omega_j
\]
for all sufficiently large $j$.

Thus
\begin{equation}
\Omega_j\overset{K}{\longrightarrow}\Omega.
\label{eq:compactness-K-final}
\end{equation}

It remains to verify that $\Omega$ belongs to the admissible class.
The uniform $C^{2,\alpha}$ bound passes to the limit:
\[
\|h\|_{C^{2,\alpha}(M)}\le A.
\]
The non-degeneracy condition is given by
\eqref{eq:compactness-limit-nondegeneracy-neighborhood}.

Moreover, the uniform geometric separation assumptions in the
definition of $\mathcal A_C(M)$, together with
\[
d_H(\partial\Omega_j,\partial\Omega)\to0,
\]
imply
\[
C\subset\Omega.
\]
Likewise, if the admissible class is defined with a fixed compact set
$K_{\rm ext}\subset M\setminus C$ satisfying
\[
K_{\rm ext}\subset M\setminus\Omega_j
\qquad\text{for every }j,
\]
then the uniform convergence of the defining functions implies
\[
K_{\rm ext}\subset M\setminus\Omega.
\]
In particular, $\Omega\neq M$.

Finally, since
\[
\nu_{\Omega_j}
=
\frac{\nabla_g h_j}
{|\nabla_g h_j|_g}
\]
on $\partial\Omega_j$, the $C^2$ convergence of the defining
functions and the uniform non-degeneracy condition imply the uniform
convergence of the corresponding normal fields. The uniform tubular
neighborhood condition is therefore stable under this convergence,
possibly after decreasing the common tubular radius. Hence $\Omega$
satisfies the same RC--GNP geometric conditions and
\[
\Omega\in\mathcal A_C(M).
\]

We have thus proved the Hausdorff convergence of the closures:
\[
d_H(\overline{\Omega_j},\overline{\Omega})
\longrightarrow0.
\]

Finally, we prove the strong $L^1$ convergence of the characteristic
functions. Let
\[
x\in M\setminus\partial\Omega.
\]
If $x\in\Omega$, then there exists a compact neighborhood
$K_x\subset\Omega$ containing $x$. By compact convergence,
$x\in\Omega_j$ for all sufficiently large $j$, and therefore
\[
\chi_{\Omega_j}(x)\longrightarrow1=\chi_\Omega(x).
\]
If
$x\in M\setminus\overline{\Omega}$, then there exists a compact
neighborhood
$K_x\subset M\setminus\overline{\Omega}$ containing $x$. Hence
$x\notin\Omega_j$ for all sufficiently large $j$, and
\[
\chi_{\Omega_j}(x)\longrightarrow0=\chi_\Omega(x).
\]
Thus
\[
\chi_{\Omega_j}(x)\longrightarrow\chi_\Omega(x)
\]
for every
$x\in M\setminus\partial\Omega$.

Since $\partial\Omega$ is a $C^{2,\alpha}$ hypersurface,
\[
dV_g(\partial\Omega)=0.
\]
Therefore
\[
\chi_{\Omega_j}\longrightarrow\chi_\Omega
\qquad\text{almost everywhere on }M.
\]
Moreover,
\[
|\chi_{\Omega_j}-\chi_\Omega|\le1,
\]
and
\[
\int_M1\,dV_g<\infty
\]
because $M$ is compact. The dominated convergence theorem gives
\[
\int_M
|\chi_{\Omega_j}-\chi_\Omega|\,dV_g
\longrightarrow0.
\]
Consequently,
\[
\chi_{\Omega_j}\longrightarrow\chi_\Omega
\qquad\text{strongly in }L^1(M,dV_g).
\]
This completes the proof.
\end{proof}

\end{document}
